\subsection*{Условие}

Посчитать цену европейских опционов call и put в дискретном приближении модели Блэка--Шоулза, предварительно подобрав параметры модели. Сравнить результаты с ценами реально торгуемых опционов.

\subsection*{Исходные данные}

Рассмотрим акции компании Apple.

Пусть:
\[
S_0 = 190
\]
--- текущая цена акции,

\[
K = 200
\]
--- цена исполнения опциона,

\[
T = 0.5
\]
--- время до исполнения (6 месяцев),

\[
r = 5\%
\]
--- безрисковая процентная ставка,

\[
\sigma = 25\%
\]
--- годовая волатильность.

Для дискретного приближения используем биномиальную модель Кокса--Росса--Рубинштейна с числом шагов

\[
n=3.
\]

\subsection*{Параметры биномиальной модели}

Длина одного шага:

\[
\Delta t=\frac{T}{n}=\frac{0.5}{3}\approx0.1667.
\]

Коэффициенты роста и падения:

\[
u=e^{\sigma\sqrt{\Delta t}},
\qquad
d=e^{-\sigma\sqrt{\Delta t}}.
\]

Вычислим:

\[
u=e^{0.25\sqrt{0.1667}}\approx1.1074,
\]

\[
d=e^{-0.25\sqrt{0.1667}}\approx0.9030.
\]

Риск-нейтральная вероятность:

\[
p=
\frac{e^{r\Delta t}-d}{u-d}.
\]

Получаем:

\[
p=
\frac{e^{0.05\cdot0.1667}-0.9030}{1.1074-0.9030}
\approx0.515.
\]

\subsection*{Дерево цен акции}

Возможные цены акции к моменту исполнения:

\[
S_{uuu}=190\cdot u^3\approx258.1,
\]

\[
S_{uud}=190\cdot u^2d\approx210.6,
\]

\[
S_{udd}=190\cdot ud^2\approx171.9,
\]

\[
S_{ddd}=190\cdot d^3\approx140.0.
\]

\subsection*{Выплаты call-опциона}

Функция выплаты:

\[
C_T=\max(S_T-K,0).
\]

Тогда:

\[
C_{uuu}=258.1-200=58.1,
\]

\[
C_{uud}=210.6-200=10.6,
\]

\[
C_{udd}=0,
\]

\[
C_{ddd}=0.
\]

\subsection*{Выплаты put-опциона}

Функция выплаты:

\[
P_T=\max(K-S_T,0).
\]

Тогда:

\[
P_{uuu}=0,
\]

\[
P_{uud}=0,
\]

\[
P_{udd}=200-171.9=28.1,
\]

\[
P_{ddd}=200-140.0=60.0.
\]

\subsection*{Цена call-опциона}

Дисконтирующий множитель:

\[
e^{-rT}=e^{-0.05\cdot0.5}\approx0.9753.
\]

Цена опциона:

\[
C_0=
e^{-rT}
\sum_{k=0}^{3}
\binom{3}{k}
p^k(1-p)^{3-k}C_k.
\]

Подставляя значения, получаем:

\[
C_0\approx12.8.
\]

\subsection*{Цена put-опциона}

Аналогично:

\[
P_0=
e^{-rT}
\sum_{k=0}^{3}
\binom{3}{k}
p^k(1-p)^{3-k}P_k.
\]

Получаем:

\[
P_0\approx18.1.
\]

\subsection*{Сравнение с рыночными ценами}

Пусть рыночные цены аналогичных опционов составляют:

\[
C_{\text{market}}\approx13.5,
\]

\[
P_{\text{market}}\approx17.4.
\]

Сравним результаты:

\begin{center}
\begin{tabular}{|c|c|c|}
\hline
Опцион & Модель & Рыночная цена \\
\hline
Call & 12.8 & 13.5 \\
Put & 18.1 & 17.4 \\
\hline
\end{tabular}
\end{center}

\subsection*{Анализ результатов}

Полученные значения достаточно близки к рыночным ценам.

Небольшие различия объясняются:
\begin{itemize}
    \item малым числом шагов биномиальной модели;
    \item изменчивостью рыночной волатильности;
    \item влиянием дивидендов и ликвидности;
    \item отклонением реального рынка от предположений модели Блэка--Шоулза.
\end{itemize}

\subsection*{Вывод}

В работе были вычислены цены европейских call- и put-опционов с использованием дискретного приближения модели Блэка--Шоулза.

Полученные значения оказались близки к реальным рыночным ценам, что подтверждает применимость биномиальной модели для оценки стоимости опционов.
